% © 2026 Ing. David Jaroš — CC BY-NC-ND 4.0
%
% This work is licensed under a Creative Commons Attribution-NonCommercial-NoDerivatives
% 4.0 International License (CC BY-NC-ND 4.0).
%
% See LICENSE.md for full license text.

% UBT-AI-PROVENANCE-BEGIN
% schema: ubt-ai-provenance/v1
% tier: C_working
% ai_assistance: disclosed
% human_review: risk-based
% editorial_responsibility: Ing. David Jaroš
% policy: ../../AI_PROVENANCE.md
% notice: Working material; exhaustive human review is not claimed.
% UBT-AI-PROVENANCE-END

\documentclass[12pt,a4paper]{article}
\usepackage[a4paper,margin=2.5cm]{geometry}
\usepackage{amsmath,amssymb,tcolorbox}
\tcbuselibrary{skins,breakable}
\newtcolorbox{statusbox}[1][]{colback=gray!8!white,colframe=black!60,arc=3pt,boxrule=1pt,title=Status Summary,#1}
\title{Vacuum graviton energy in UBT and cosmological constant scope note}
\author{Ing.\ David Jaroš}
\date{2026-05-12}
\begin{document}
\maketitle

Using the graviton mode framework from
\texttt{research\_tracks/quantum\_ubt/graviton\_quantisation.tex}, define
\[
E_{\mathrm{vac}}^\zeta = \mu^{2s}\,\frac12\,\zeta_D\!\left(s-\frac12\right)\Big|_{s=0},
\qquad
\zeta_D(s)=\sum_{\ell=2}^{\infty}(2\ell+1)\,\omega_\ell^{-2s}.
\]
A simple asymptotic model $\omega_\ell\sim (\ell+\tfrac12)/r_s$ gives the standard
large vacuum-energy hierarchy (Planck-scale mismatch), i.e. no automatic
120-order suppression from this first pass.

\begin{statusbox}
\textbf{Hlavní výsledek}: Formální $\zeta$-regularizační rámec je definován; první-pass odhad reprodukuje standardní kosmologický problém vakuové energie.\\
\textbf{Klasifikace}: [STD framework + OPEN mechanism].\\
\textbf{Dopad na teorii}: V této fázi není identifikován nový UBT supresní mechanismus kosmologické konstanty.\\
\textbf{Příští krok}: Nahradit asymptotické $\omega_\ell$ numerickým spektrem z RW/Zerilli a otestovat, zda UBT struktura mění UV asymptotiku.\\
\textbf{Alpha je odvozena}: NE (unconditional)
\end{statusbox}

\end{document}
